% Generated draft: requires mathematical and source review.
\documentclass[11pt]{article}
\usepackage[margin=25mm]{geometry}
\usepackage{fontspec}
\setmainfont{Latin Modern Roman}
\usepackage{amsmath,amssymb,mathtools,unicode-math}
\setmathfont{Latin Modern Math}
\usepackage{enumitem,needspace,xurl,hyperref}
\setlength{\parindent}{0pt}
\setlength{\parskip}{5pt}
\setlength{\emergencystretch}{3em}
\newcommand{\N}{\mathbb{N}}
\newcommand{\Z}{\mathbb{Z}}
\newcommand{\Q}{\mathbb{Q}}
\newcommand{\R}{\mathbb{R}}
\newcommand{\C}{\mathbb{C}}
\newcommand{\F}{\mathbb{F}}
\newcommand{\A}{\mathbb{A}}
\newcommand{\PP}{\mathbb{P}}
\newcommand{\E}{\mathbb{E}}
\newcommand{\eps}{\varepsilon}
\newcommand{\dd}{\,\mathrm{d}}
\begin{document}
\section*{UnsolvedMath ID: 30002029}
\textbf{Source problem code:} catalogue entry
\subsection{Definitions and mathematical statement}
For a positive integer \(r\), a finite simple \(r\)-uniform hypergraph is a pair \(H=(V(H),E(H))\), where \(V(H)\) is finite and \(E(H)\subseteq\binom{V(H)}{r}\). A copy of a hypergraph \(G\) in \(H\) means an injective map \(f:V(G)\to V(H)\) such that \(f(e)\in E(H)\) for every \(e\in E(G)\); the copy need not be induced. For a finite family \(F\) of finite \(r\)-uniform hypergraphs, call \(H\) \(F\)-free if it contains no copy of any member of \(F\).

In the usual nondegenerate setting, where every member of \(F\) has at least one edge, define
\[
\operatorname{ex}(n,F)=\max\bigl\{|E(H)|:|V(H)|=n,\ H\text{ is }F\text{-free}\bigr\},
\qquad
\pi(F)=\lim_{n\to\infty}\frac{\operatorname{ex}(n,F)}{\binom{n}{r}}.
\]
The limit is taken over integers \(n\ge r\). These definitions use the usual simple, non-induced hypergraph convention. The supplied record does not explicitly specify these conventions, the permitted range of \(r\), or the treatment of forbidden hypergraphs with no edges, for which the displayed maximum can have an empty feasible set.

\textbf{Decision problem.} Given a finite family \(F\) of \(r\)-uniform hypergraphs and a rational number \(\alpha\), is it decidable whether
\[
\pi(F)\leq\alpha?
\]
Here decidability means that an algorithm, given finite encodings of the input hypergraphs and of \(\alpha\), always halts and answers yes exactly when the displayed inequality holds, including equality, and answers no otherwise.

The source leaves the role of \(r\) unresolved. A uniform formulation asks for one algorithm taking \((r,F,\alpha)\) as input. A fixed-uniformity formulation asks, for each permitted fixed \(r\), whether there exists an algorithm taking \((F,\alpha)\) as input. These have different quantifier orders; the supplied wording does not select between them.
\subsection{Short English statement}
Can an algorithm always decide whether the limiting maximum edge density of hypergraphs avoiding a given finite family is at most a given rational number? Equality is included. The intended quantification over the uniformity \(r\) and conventions for degenerate inputs require clarification. This is a statement-only draft requiring human mathematical review.
\subsection{Sources}
\begin{itemize}
\item Supplied local registry record: permanent UnsolvedMath ID \(30002029\), source problem code \texttt{catalogue entry}, title ``Decidability of Hypergraph Tur\'an Densities,'' in the Oberwolfach Reports Open Problems collection. No individual proposer is identified in the supplied record.
\item Supplied source citation: \emph{Mini-Workshop: Hypergraph Tur\'an Problem (2013)}, Oberwolfach Reports, pp.\ 1149--1182, DOI \texttt{10.4171/owr/2012/19}. \url{https://doi.org/10.4171/owr/2012/19}. The registry separately gives proposed year 2012 and an \texttt{owr-year:2012} tag; the relationship to the citation's 2013 label is unresolved.
\item Supplied literature source, ``Oberwolfach report'': \url{https://ems.press/content/serial-article-files/46390?nt=1}.
\item Supplied literature source, ``Pikhurko on possible Turan densities'': \url{https://doi.org/10.1007/s11856-014-0031-5}.
\item Supplied literature source, ``Pikhurko on large algebraic degree'': \url{https://wrap.warwick.ac.uk/174830/1/WRAP-hypergraph-turan-densities-arbitrarily-large-algebraic-degree-Pikhurko-2023.pdf}.
\item Literature assessment dated exactly \texttt{2026-08-21}: the supplied curation labels the problem open and reports that no algorithm deciding this inequality for arbitrary finite hypergraph families has been proved, and that subsequent work on possible Tur\'an densities and high algebraic degree does not resolve decidability. This is supplied dated curation derived from the OWR record, not an independently verified literature assessment. The source requests independent verification of time-sensitive status and cited claims.
\item The registry labels the statement \texttt{corrected\_verified} and reports a manual semantic, typographic, and LaTeX review of a meaning-preserving cleanup. This supplied editorial label does not constitute independent mathematical review of this draft.
\end{itemize}
\end{document}
